\IfFileExists{stacks-project.cls}{%
\documentclass{stacks-project}
}{%
\documentclass{amsart}
}

% For dealing with references we use the comment environment
\usepackage{verbatim}
\newenvironment{reference}{\comment}{\endcomment}
%\newenvironment{reference}{}{}
\newenvironment{slogan}{\comment}{\endcomment}
\newenvironment{history}{\comment}{\endcomment}

% For commutative diagrams we use Xy-pic
\usepackage[all]{xy}

% We use 2cell for 2-commutative diagrams.
\xyoption{2cell}
\UseAllTwocells

% We use multicol for the list of chapters between chapters
\usepackage{multicol}

% This is generally recommended for better output
\usepackage{lmodern}
\usepackage[T1]{fontenc}

% For cross-file-references
\usepackage{xr-hyper}

% Package for hypertext links:
\usepackage{hyperref}

% For any local file, say "hello.tex" you want to link to please
% use \externaldocument[hello-]{hello}
\externaldocument[introduction-]{introduction}
\externaldocument[conventions-]{conventions}
\externaldocument[sets-]{sets}
\externaldocument[categories-]{categories}
\externaldocument[topology-]{topology}
\externaldocument[sheaves-]{sheaves}
\externaldocument[sites-]{sites}
\externaldocument[stacks-]{stacks}
\externaldocument[fields-]{fields}
\externaldocument[algebra-]{algebra}
\externaldocument[brauer-]{brauer}
\externaldocument[homology-]{homology}
\externaldocument[derived-]{derived}
\externaldocument[simplicial-]{simplicial}
\externaldocument[more-algebra-]{more-algebra}
\externaldocument[smoothing-]{smoothing}
\externaldocument[modules-]{modules}
\externaldocument[sites-modules-]{sites-modules}
\externaldocument[injectives-]{injectives}
\externaldocument[cohomology-]{cohomology}
\externaldocument[sites-cohomology-]{sites-cohomology}
\externaldocument[dga-]{dga}
\externaldocument[dpa-]{dpa}
\externaldocument[sdga-]{sdga}
\externaldocument[hypercovering-]{hypercovering}
\externaldocument[schemes-]{schemes}
\externaldocument[constructions-]{constructions}
\externaldocument[properties-]{properties}
\externaldocument[morphisms-]{morphisms}
\externaldocument[coherent-]{coherent}
\externaldocument[divisors-]{divisors}
\externaldocument[limits-]{limits}
\externaldocument[varieties-]{varieties}
\externaldocument[topologies-]{topologies}
\externaldocument[descent-]{descent}
\externaldocument[perfect-]{perfect}
\externaldocument[more-morphisms-]{more-morphisms}
\externaldocument[flat-]{flat}
\externaldocument[groupoids-]{groupoids}
\externaldocument[more-groupoids-]{more-groupoids}
\externaldocument[etale-]{etale}
\externaldocument[chow-]{chow}
\externaldocument[intersection-]{intersection}
\externaldocument[pic-]{pic}
\externaldocument[weil-]{weil}
\externaldocument[adequate-]{adequate}
\externaldocument[dualizing-]{dualizing}
\externaldocument[duality-]{duality}
\externaldocument[discriminant-]{discriminant}
\externaldocument[derham-]{derham}
\externaldocument[local-cohomology-]{local-cohomology}
\externaldocument[algebraization-]{algebraization}
\externaldocument[curves-]{curves}
\externaldocument[resolve-]{resolve}
\externaldocument[models-]{models}
\externaldocument[functors-]{functors}
\externaldocument[equiv-]{equiv}
\externaldocument[pione-]{pione}
\externaldocument[etale-cohomology-]{etale-cohomology}
\externaldocument[proetale-]{proetale}
\externaldocument[relative-cycles-]{relative-cycles}
\externaldocument[more-etale-]{more-etale}
\externaldocument[trace-]{trace}
\externaldocument[crystalline-]{crystalline}
\externaldocument[spaces-]{spaces}
\externaldocument[spaces-properties-]{spaces-properties}
\externaldocument[spaces-morphisms-]{spaces-morphisms}
\externaldocument[decent-spaces-]{decent-spaces}
\externaldocument[spaces-cohomology-]{spaces-cohomology}
\externaldocument[spaces-limits-]{spaces-limits}
\externaldocument[spaces-divisors-]{spaces-divisors}
\externaldocument[spaces-over-fields-]{spaces-over-fields}
\externaldocument[spaces-topologies-]{spaces-topologies}
\externaldocument[spaces-descent-]{spaces-descent}
\externaldocument[spaces-perfect-]{spaces-perfect}
\externaldocument[spaces-more-morphisms-]{spaces-more-morphisms}
\externaldocument[spaces-flat-]{spaces-flat}
\externaldocument[spaces-groupoids-]{spaces-groupoids}
\externaldocument[spaces-more-groupoids-]{spaces-more-groupoids}
\externaldocument[bootstrap-]{bootstrap}
\externaldocument[spaces-pushouts-]{spaces-pushouts}
\externaldocument[spaces-chow-]{spaces-chow}
\externaldocument[groupoids-quotients-]{groupoids-quotients}
\externaldocument[spaces-more-cohomology-]{spaces-more-cohomology}
\externaldocument[spaces-simplicial-]{spaces-simplicial}
\externaldocument[spaces-duality-]{spaces-duality}
\externaldocument[formal-spaces-]{formal-spaces}
\externaldocument[restricted-]{restricted}
\externaldocument[spaces-resolve-]{spaces-resolve}
\externaldocument[formal-defos-]{formal-defos}
\externaldocument[defos-]{defos}
\externaldocument[cotangent-]{cotangent}
\externaldocument[examples-defos-]{examples-defos}
\externaldocument[algebraic-]{algebraic}
\externaldocument[examples-stacks-]{examples-stacks}
\externaldocument[stacks-sheaves-]{stacks-sheaves}
\externaldocument[criteria-]{criteria}
\externaldocument[artin-]{artin}
\externaldocument[quot-]{quot}
\externaldocument[stacks-properties-]{stacks-properties}
\externaldocument[stacks-morphisms-]{stacks-morphisms}
\externaldocument[stacks-limits-]{stacks-limits}
\externaldocument[stacks-cohomology-]{stacks-cohomology}
\externaldocument[stacks-perfect-]{stacks-perfect}
\externaldocument[stacks-introduction-]{stacks-introduction}
\externaldocument[stacks-more-morphisms-]{stacks-more-morphisms}
\externaldocument[stacks-geometry-]{stacks-geometry}
\externaldocument[moduli-]{moduli}
\externaldocument[moduli-curves-]{moduli-curves}
\externaldocument[examples-]{examples}
\externaldocument[exercises-]{exercises}
\externaldocument[guide-]{guide}
\externaldocument[desirables-]{desirables}
\externaldocument[coding-]{coding}
\externaldocument[obsolete-]{obsolete}
\externaldocument[fdl-]{fdl}
\externaldocument[index-]{index}

% Theorem environments.
%
\theoremstyle{plain}
\newtheorem{theorem}[subsection]{Theorem}
\newtheorem{proposition}[subsection]{Proposition}
\newtheorem{lemma}[subsection]{Lemma}

\theoremstyle{definition}
\newtheorem{definition}[subsection]{Definition}
\newtheorem{example}[subsection]{Example}
\newtheorem{exercise}[subsection]{Exercise}
\newtheorem{situation}[subsection]{Situation}

\theoremstyle{remark}
\newtheorem{remark}[subsection]{Remark}
\newtheorem{remarks}[subsection]{Remarks}

\numberwithin{equation}{subsection}

% Macros
%
\def\lim{\mathop{\mathrm{lim}}\nolimits}
\def\colim{\mathop{\mathrm{colim}}\nolimits}
\def\Spec{\mathop{\mathrm{Spec}}}
\def\Hom{\mathop{\mathrm{Hom}}\nolimits}
\def\Ext{\mathop{\mathrm{Ext}}\nolimits}
\def\SheafHom{\mathop{\mathcal{H}\!\mathit{om}}\nolimits}
\def\SheafExt{\mathop{\mathcal{E}\!\mathit{xt}}\nolimits}
\def\Sch{\mathit{Sch}}
\def\Mor{\mathop{\mathrm{Mor}}\nolimits}
\def\Ob{\mathop{\mathrm{Ob}}\nolimits}
\def\Sh{\mathop{\mathit{Sh}}\nolimits}
\def\NL{\mathop{N\!L}\nolimits}
\def\CH{\mathop{\mathrm{CH}}\nolimits}
\def\proetale{{pro\text{-}\acute{e}tale}}
\def\etale{{\acute{e}tale}}
\def\QCoh{\mathit{QCoh}}
\def\Ker{\mathop{\mathrm{Ker}}}
\def\Im{\mathop{\mathrm{Im}}}
\def\Coker{\mathop{\mathrm{Coker}}}
\def\Coim{\mathop{\mathrm{Coim}}}

% Boxtimes
%
\DeclareMathSymbol{\boxtimes}{\mathbin}{AMSa}{"02}

%
% Macros for moduli stacks/spaces
%
\def\QCohstack{\mathcal{QC}\!\mathit{oh}}
\def\Cohstack{\mathcal{C}\!\mathit{oh}}
\def\Spacesstack{\mathcal{S}\!\mathit{paces}}
\def\Quotfunctor{\mathrm{Quot}}
\def\Hilbfunctor{\mathrm{Hilb}}
\def\Curvesstack{\mathcal{C}\!\mathit{urves}}
\def\Polarizedstack{\mathcal{P}\!\mathit{olarized}}
\def\Complexesstack{\mathcal{C}\!\mathit{omplexes}}
% \Pic is the operator that assigns to X its picard group, usage \Pic(X)
% \Picardstack_{X/B} denotes the Picard stack of X over B
% \Picardfunctor_{X/B} denotes the Picard functor of X over B
\def\Pic{\mathop{\mathrm{Pic}}\nolimits}
\def\Picardstack{\mathcal{P}\!\mathit{ic}}
\def\Picardfunctor{\mathrm{Pic}}
\def\Deformationcategory{\mathcal{D}\!\mathit{ef}}
\usepackage{cleveref}
\crefname{originaldefinition}{original definition}{original definitions}
\Crefname{originaldefinition}{Original definition}{Original definitions}
\makeatletter
\newlabel{spaces-more-morphisms-definition-relatively-perfect}{{0DKN}{source}{}{sourcepointer-1}{}}
\newlabel{spaces-more-morphisms-definition-relatively-perfect@cref}{{[originaldefinition][0][]0DKN}{[1][0][]source}{}{}{}}
\makeatother
\crefname{originallemma}{original lemma}{original lemmas}
\Crefname{originallemma}{Original lemma}{Original lemmas}
\makeatletter
\newlabel{sites-cohomology-lemma-pseudo-coherent-pullback}{{08H4}{source}{}{sourcepointer-2}{}}
\newlabel{sites-cohomology-lemma-pseudo-coherent-pullback@cref}{{[originallemma][0][]08H4}{[1][0][]source}{}{}{}}
\makeatother
\crefname{originallemma}{original lemma}{original lemmas}
\Crefname{originallemma}{Original lemma}{Original lemmas}
\makeatletter
\newlabel{spaces-more-morphisms-lemma-base-change-relatively-perfect}{{0DKT}{source}{}{sourcepointer-3}{}}
\newlabel{spaces-more-morphisms-lemma-base-change-relatively-perfect@cref}{{[originallemma][0][]0DKT}{[1][0][]source}{}{}{}}
\makeatother
\crefname{originallemma}{original lemma}{original lemmas}
\Crefname{originallemma}{Original lemma}{Original lemmas}
\makeatletter
\newlabel{perfect-lemma-affine-morphism-and-hom-out-of-perfect}{{0AA7}{source}{}{sourcepointer-4}{}}
\newlabel{perfect-lemma-affine-morphism-and-hom-out-of-perfect@cref}{{[originallemma][0][]0AA7}{[1][0][]source}{}{}{}}
\makeatother
\crefname{originallemma}{original lemma}{original lemmas}
\Crefname{originallemma}{Original lemma}{Original lemmas}
\makeatletter
\newlabel{spaces-more-morphisms-lemma-compute-ext-rel-perfect}{{0DKY}{source}{}{sourcepointer-5}{}}
\newlabel{spaces-more-morphisms-lemma-compute-ext-rel-perfect@cref}{{[originallemma][0][]0DKY}{[1][0][]source}{}{}{}}
\makeatother
\crefname{originallemma}{original lemma}{original lemmas}
\Crefname{originallemma}{Original lemma}{Original lemmas}
\makeatletter
\newlabel{perfect-lemma-jump-loci}{{0BDI}{source}{}{sourcepointer-6}{}}
\newlabel{perfect-lemma-jump-loci@cref}{{[originallemma][0][]0BDI}{[1][0][]source}{}{}{}}
\makeatother
\crefname{originallemma}{original lemma}{original lemmas}
\Crefname{originallemma}{Original lemma}{Original lemmas}
\makeatletter
\newlabel{sites-cohomology-lemma-section-RHom-over-U}{{08JA}{source}{}{sourcepointer-7}{}}
\newlabel{sites-cohomology-lemma-section-RHom-over-U@cref}{{[originallemma][0][]08JA}{[1][0][]source}{}{}{}}
\makeatother
\crefname{originallemma}{original lemma}{original lemmas}
\Crefname{originallemma}{Original lemma}{Original lemmas}
\makeatletter
\newlabel{sites-cohomology-lemma-Leray-unbounded}{{0D6H}{source}{}{sourcepointer-8}{}}
\newlabel{sites-cohomology-lemma-Leray-unbounded@cref}{{[originallemma][0][]0D6H}{[1][0][]source}{}{}{}}
\makeatother
\crefname{originallemma}{original lemma}{original lemmas}
\Crefname{originallemma}{Original lemma}{Original lemmas}
\makeatletter
\newlabel{derived-lemma-negative-vanishing}{{05TC}{source}{}{sourcepointer-9}{}}
\newlabel{derived-lemma-negative-vanishing@cref}{{[originallemma][0][]05TC}{[1][0][]source}{}{}{}}
\makeatother
\crefname{originallemma}{original lemma}{original lemmas}
\Crefname{originallemma}{Original lemma}{Original lemmas}
\makeatletter
\newlabel{more-algebra-lemma-lift-pseudo-coherent-from-residue-field}{{0BCC}{source}{}{sourcepointer-10}{}}
\newlabel{more-algebra-lemma-lift-pseudo-coherent-from-residue-field@cref}{{[originallemma][0][]0BCC}{[1][0][]source}{}{}{}}
\makeatother
\crefname{originallemma}{original lemma}{original lemmas}
\Crefname{originallemma}{Original lemma}{Original lemmas}
\makeatletter
\newlabel{spaces-simplicial-lemma-fppf-neg-ext-zero-hom}{{0DL9}{source}{}{sourcepointer-11}{}}
\newlabel{spaces-simplicial-lemma-fppf-neg-ext-zero-hom@cref}{{[originallemma][0][]0DL9}{[1][0][]source}{}{}{}}
\makeatother
\crefname{originallemma}{original lemma}{original lemmas}
\Crefname{originallemma}{Original lemma}{Original lemmas}
\makeatletter
\newlabel{bootstrap-lemma-locally-algebraic-space-finite-type}{{04U0}{source}{}{sourcepointer-12}{}}
\newlabel{bootstrap-lemma-locally-algebraic-space-finite-type@cref}{{[originallemma][0][]04U0}{[1][0][]source}{}{}{}}
\makeatother
\crefname{originallemma}{original lemma}{original lemmas}
\Crefname{originallemma}{Original lemma}{Original lemmas}
\makeatletter
\newlabel{spaces-morphisms-lemma-affine-local}{{03WG}{source}{}{sourcepointer-13}{}}
\newlabel{spaces-morphisms-lemma-affine-local@cref}{{[originallemma][0][]03WG}{[1][0][]source}{}{}{}}
\makeatother
\crefname{originallemma}{original lemma}{original lemmas}
\Crefname{originallemma}{Original lemma}{Original lemmas}
\makeatletter
\newlabel{spaces-morphisms-lemma-finite-presentation-local}{{0410}{source}{}{sourcepointer-14}{}}
\newlabel{spaces-morphisms-lemma-finite-presentation-local@cref}{{[originallemma][0][]0410}{[1][0][]source}{}{}{}}
\makeatother
\crefname{originallemma}{original lemma}{original lemmas}
\Crefname{originallemma}{Original lemma}{Original lemmas}
\makeatletter
\newlabel{more-algebra-lemma-pseudo-coherent}{{064U}{source}{}{sourcepointer-15}{}}
\newlabel{more-algebra-lemma-pseudo-coherent@cref}{{[originallemma][0][]064U}{[1][0][]source}{}{}{}}
\makeatother
\title{Open locus of negative Ext vanishing for a pseudo-coherent complex against a relatively perfect one, with affine finitely presented Hom space}
\author{The Stacks Project Authors}
\date{stacks-project:quot.tex:local-snapshot-20261012; GNU FDL 1.2 or later}
\begin{document}
\maketitle
\setcounter{subsection}{1}
\noindent\textit{Editorial note.} All delivered mathematical text is the author's own, copied from the pinned Stacks Project chapter file quot.tex; no statement or proof step was rewritten, reordered or removed. The counted span is the author environment and its own proof environment. References inside the span are handed to the assembler as navigation stubs that carry the canonical tag of the cited tag only; whether each of them is load-bearing is recorded on the stub itself.\par\medskip
\noindent\textit{Allowed standard prerequisites (editorial).}\begin{itemize}
\item Definition spaces-more-morphisms-definition-relatively-perfect (canonical tag 0DKN). cited at statement lines 5159-5160. E is Y-perfect.
\item Lemma sites-cohomology-lemma-pseudo-coherent-pullback (canonical tag 08H4). cited at proof lines 5185-5187. K\_V stays pseudo-coherent under base change.
\item Lemma spaces-more-morphisms-lemma-base-change-relatively-perfect (canonical tag 0DKT). cited at proof lines 5188-5189. E\_V stays V-perfect under base change.
\item Lemma perfect-lemma-affine-morphism-and-hom-out-of-perfect (canonical tag 0AA7). cited at proof lines 5193-5198. Ext over fibres is stable under field extension.
\item Lemma spaces-more-morphisms-lemma-compute-ext-rel-perfect (canonical tag 0DKY). cited at proof lines 5204-5206. L universally computes Rf\_*RSheafHom(K, E) over affine Y.
\item Lemma perfect-lemma-jump-loci (canonical tag 0BDI). cited at proof line 5221. The vanishing locus of H\textasciicircum{}i of a pseudo-coherent complex on fibres is open.
\item Lemma sites-cohomology-lemma-section-RHom-over-U (canonical tag 08JA). cited at proof lines 5245-5246. Ext\textasciicircum{}i as H\textasciicircum{}i(X\_V, RSheafHom).
\item Lemma sites-cohomology-lemma-Leray-unbounded (canonical tag 0D6H). cited at proof lines 5246-5247. H\textasciicircum{}i(X\_V, -) = H\textasciicircum{}i(V, Rf\_\{V,*\}-).
\item Lemma derived-lemma-negative-vanishing (canonical tag 05TC). cited at proof lines 5249-5250. Vanishing cohomology sheaves in degrees < 0 give vanishing H\textasciicircum{}i, i < 0.
\item Lemma more-algebra-lemma-lift-pseudo-coherent-from-residue-field (canonical tag 0BCC). cited at proof lines 5260-5261. Fibre-wise vanishing lifts to stalks for pseudo-coherent complexes.
\item Lemma spaces-simplicial-lemma-fppf-neg-ext-zero-hom (canonical tag 0DL9). cited at proof lines 5274 and 5283 (footnote). F is an fppf sheaf via the Cech hypercovering argument in the footnote.
\item Lemma bootstrap-lemma-locally-algebraic-space-finite-type (canonical tag 04U0). cited at proof line 5287. Representability is etale local.
\item Lemma spaces-morphisms-lemma-affine-local (canonical tag 03WG). cited at proof line 5289. Affineness is etale local.
\item Lemma spaces-morphisms-lemma-finite-presentation-local (canonical tag 0410). cited at proof lines 5289-5290. Finite presentation is etale local.
\item Lemma more-algebra-lemma-pseudo-coherent (canonical tag 064U). cited at proof lines 5300-5301. L is represented by a two-term complex of finite free modules in degrees 1, 0.
\end{itemize}\medskip
\ifdefined\hypertarget\hypertarget{sourcepointer-1}{}\fi
\noindent\textit{Source reference.} Original definition 0DKN: sources/stacks-chow/tags.txt, source lines 17018--17018. Cited at statement lines 5159-5160: E is Y-perfect per this definition; load-bearing.\par\medskip
\ifdefined\hypertarget\hypertarget{sourcepointer-2}{}\fi
\noindent\textit{Source reference.} Original lemma 08H4: sources/stacks-chow/tags.txt, source lines 10767--10767. Cited at proof lines 5185-5187: K\_V is pseudo-coherent for any V -> Y; load-bearing (first observation).\par\medskip
\ifdefined\hypertarget\hypertarget{sourcepointer-3}{}\fi
\noindent\textit{Source reference.} Original lemma 0DKT: sources/stacks-chow/tags.txt, source lines 17023--17023. Cited at proof lines 5188-5189: E\_V is V-perfect; load-bearing (first observation).\par\medskip
\ifdefined\hypertarget\hypertarget{sourcepointer-4}{}\fi
\noindent\textit{Source reference.} Original lemma 0AA7: sources/stacks-chow/tags.txt, source lines 12975--12975. Cited at proof lines 5197-5198: Ext\textasciicircum{}i over the fibre tensor k' equals Ext\textasciicircum{}i over k; load-bearing (well-definedness on |Y|).\par\medskip
\ifdefined\hypertarget\hypertarget{sourcepointer-5}{}\fi
\noindent\textit{Source reference.} Original lemma 0DKY: sources/stacks-chow/tags.txt, source lines 17028--17028. Cited at proof lines 5204-5206: supplies L universally computing Rf\_*RSheafHom(K, E) over affine Y; load-bearing.\par\medskip
\ifdefined\hypertarget\hypertarget{sourcepointer-6}{}\fi
\noindent\textit{Source reference.} Original lemma 0BDI: sources/stacks-chow/tags.txt, source lines 14316--14316. Cited at proof line 5221: the locus where H\textasciicircum{}i(L x\textasciicircum{}L kappa(y)) = 0 for i > 0 is open; load-bearing for part (1).\par\medskip
\ifdefined\hypertarget\hypertarget{sourcepointer-7}{}\fi
\noindent\textit{Source reference.} Original lemma 08JA: sources/stacks-chow/tags.txt, source lines 10843--10843. Cited at proof lines 5245-5246: first equality expressing Ext\textasciicircum{}i over X\_V as H\textasciicircum{}i(X\_V, RSheafHom); load-bearing for part (2).\par\medskip
\ifdefined\hypertarget\hypertarget{sourcepointer-8}{}\fi
\noindent\textit{Source reference.} Original lemma 0D6H: sources/stacks-chow/tags.txt, source lines 16522--16522. Cited at proof lines 5246-5247: second equality H\textasciicircum{}i(X\_V, -) = H\textasciicircum{}i(V, Rf\_\{V,*\}-); load-bearing for part (2).\par\medskip
\ifdefined\hypertarget\hypertarget{sourcepointer-9}{}\fi
\noindent\textit{Source reference.} Original lemma 05TC: sources/stacks-chow/tags.txt, source lines 7476--7476. Cited at proof lines 5249-5250: vanishing cohomology sheaves in degrees < 0 kills H\textasciicircum{}i for i < 0; load-bearing for part (2).\par\medskip
\ifdefined\hypertarget\hypertarget{sourcepointer-10}{}\fi
\noindent\textit{Source reference.} Original lemma 0BCC: sources/stacks-chow/tags.txt, source lines 14275--14275. Cited at proof lines 5260-5261: W = Y gives H\textasciicircum{}i(L) = 0 for i > 0 on stalks; load-bearing for the claim.\par\medskip
\ifdefined\hypertarget\hypertarget{sourcepointer-11}{}\fi
\noindent\textit{Source reference.} Original lemma 0DL9: sources/stacks-chow/tags.txt, source lines 17039--17039. Cited at proof lines 5274 and 5283 (including the footnote hypercovering argument): F is an fppf sheaf; load-bearing for part (3).\par\medskip
\ifdefined\hypertarget\hypertarget{sourcepointer-12}{}\fi
\noindent\textit{Source reference.} Original lemma 04U0: sources/stacks-chow/tags.txt, source lines 6273--6273. Cited at proof line 5287: checking representability etale locally on Y; load-bearing for part (3).\par\medskip
\ifdefined\hypertarget\hypertarget{sourcepointer-13}{}\fi
\noindent\textit{Source reference.} Original lemma 03WG: sources/stacks-chow/tags.txt, source lines 5111--5111. Cited at proof line 5289: affineness of F -> Y checked etale locally; load-bearing for part (3).\par\medskip
\ifdefined\hypertarget\hypertarget{sourcepointer-14}{}\fi
\noindent\textit{Source reference.} Original lemma 0410: sources/stacks-chow/tags.txt, source lines 5275--5275. Cited at proof lines 5289-5290: finite presentation of F -> Y checked etale locally; load-bearing for part (3).\par\medskip
\ifdefined\hypertarget\hypertarget{sourcepointer-15}{}\fi
\noindent\textit{Source reference.} Original lemma 064U: sources/stacks-chow/tags.txt, source lines 7878--7878. Cited at proof lines 5300-5301: L represented by O\textasciicircum{}m\_1 -> O\textasciicircum{}m\_0 in degrees 1, 0; load-bearing for the fibre-product presentation of F.\par\medskip
\begin{lemma}
\label{lemma-complexes-open-neg-exts-vanishing}
Let $S$ be a scheme.
Let $f : X \to Y$ be a morphism of algebraic spaces over $S$.
Assume $f$ is proper, flat, and of finite presentation.
Let $K, E \in D(\mathcal{O}_X)$. Assume $K$ is pseudo-coherent
and $E$ is $Y$-perfect (More on Morphisms of Spaces, Definition
\ref{spaces-more-morphisms-definition-relatively-perfect}).
For a field $k$ and a morphism $y : \Spec(k) \to Y$ denote $K_y$, $E_y$
the pullback to the fibre $X_y$.
\begin{enumerate}
\item There is an open $W \subset Y$ characterized by the property
$$
y \in |W|
\Leftrightarrow
\Ext^i_{\mathcal{O}_{X_y}}(K_y, E_y) = 0
\text{ for }i < 0.
$$
\item For any morphism $V \to Y$ factoring through $W$ we have
$$
\Ext^i_{\mathcal{O}_{X_V}}(K_V, E_V) = 0
\quad\text{for}\quad i < 0
$$
where $X_V$ is the base change of $X$ and $K_V$ and $E_V$
are the derived pullbacks of $K$ and $E$ to $X_V$.
\item The functor $V \mapsto \Hom_{\mathcal{O}_{X_V}}(K_V, E_V)$
is a sheaf on $(\textit{Spaces}/W)_{fppf}$ representable by an
algebraic space affine and of finite presentation over $W$.
\end{enumerate}
\end{lemma}

\begin{proof}
For any morphism $V \to Y$ the complex $K_V$ is pseudo-coherent
(Cohomology on Sites, Lemma
\ref{sites-cohomology-lemma-pseudo-coherent-pullback})
and $E_V$ is $V$-perfect (More on Morphisms of Spaces, Lemma
\ref{spaces-more-morphisms-lemma-base-change-relatively-perfect}).
Another observation is that given $y : \Spec(k) \to Y$
and a field extension $k'/k$ with $y' : \Spec(k') \to Y$
the induced morphism, we have
$$
\Ext^i_{\mathcal{O}_{X_{y'}}}(K_{y'}, E_{y'}) =
\Ext^i_{\mathcal{O}_{X_y}}(K_y, E_y) \otimes_k k'
$$
by Derived Categories of Schemes, Lemma
\ref{perfect-lemma-affine-morphism-and-hom-out-of-perfect}.
Thus the vanishing in (1) is really a property of the induced
point $y \in |Y|$.
We will use these two observations without further mention in the proof.

\medskip\noindent
Assume first $Y$ is an affine scheme. Then we may apply
More on Morphisms of Spaces, Lemma
\ref{spaces-more-morphisms-lemma-compute-ext-rel-perfect}
and find a pseudo-coherent $L \in D(\mathcal{O}_Y)$ which
``universally computes'' $Rf_*R\SheafHom(K, E)$ in the sense
described in that lemma. Unwinding the definitions, we obtain
for a point $y \in Y$ the equality
$$
\Ext^i_{\kappa(y)}(L \otimes_{\mathcal{O}_Y}^\mathbf{L} \kappa(y),
\kappa(y)) = \Ext^i_{\mathcal{O}_{X_y}}(K_y, E_y)
$$
We conclude that
$$
H^i(L \otimes_{\mathcal{O}_Y}^\mathbf{L} \kappa(y)) = 0
\text{ for } i > 0 \Leftrightarrow
\Ext^i_{\mathcal{O}_{X_y}}(K_y, E_y) = 0 \text{ for }i < 0.
$$
By Derived Categories of Schemes, Lemma \ref{perfect-lemma-jump-loci}
the set $W$ of $y \in Y$ where this happens defines an open of $Y$.
This open $W$ then satisfies the requirement in (1) for all morphisms
from spectra of fields, by the ``universality'' of $L$.

\medskip\noindent
Let's go back to $Y$ a general algebraic space.
Choose an \'etale covering $\{V_i \to Y\}$ by affine schemes $V_i$.
Then we see that the subset $W \subset |Y|$ pulls back to the corresponding
subset $W_i \subset |V_i|$ for $X_{V_i}$, $K_{V_i}$, $E_{V_i}$.
By the previous paragraph we find that $W_i$ is open, hence $W$ is open.
This proves (1) in general. Moreover, parts (2) and (3) are entirely formulated
in terms of the category $\textit{Spaces}/W$ and the restrictions
$X_W$, $K_W$, $E_W$. This reduces us to the case $W = Y$.

\medskip\noindent
Assume $W = Y$. We claim that for any algebraic space $V$ over $Y$
we have $Rf_{V, *}R\SheafHom(K_V, E_V)$ has vanishing cohomology
sheaves in degrees $< 0$. This will prove (2) because
$$
\Ext^i_{\mathcal{O}_{X_V}}(K_V, E_V) =
H^i(X_V, R\SheafHom(K_V, E_V)) = 
H^i(V, Rf_{V, *}R\SheafHom(K_V, E_V))
$$
by Cohomology on Sites, Lemmas
\ref{sites-cohomology-lemma-section-RHom-over-U} and
\ref{sites-cohomology-lemma-Leray-unbounded}
and the vanishing of the cohomology sheaves implies the
cohomology group $H^i$ is zero for $i < 0$ by
Derived Categories, Lemma \ref{derived-lemma-negative-vanishing}.

\medskip\noindent
To prove the claim, we may work \'etale locally on $V$.
In particular, we may assume $Y$ is affine and $W = Y$.
Let $L \in D(\mathcal{O}_Y)$ be as in the second paragraph of the proof.
For an algebraic space $V$ over $Y$ denote $L_V$ the derived pullback of
$L$ to $V$. (An important feature we will use is that $L$ ``works'' for all
algebraic spaces $V$ over $Y$ and not just affine $V$.)
As $W = Y$ we have $H^i(L) = 0$ for $i > 0$
(use More on Algebra, Lemma
\ref{more-algebra-lemma-lift-pseudo-coherent-from-residue-field}
to go from fibres to stalks). Hence $H^i(L_V) = 0$ for $i > 0$.
The property defining $L$ is that
$$
Rf_{V, *}R\SheafHom(K_V, E_V) = R\SheafHom(L_V, \mathcal{O}_V)
$$
Since $L_V$ sits in degrees $\leq 0$, we conclude that
$R\SheafHom(L_V, \mathcal{O}_V)$ sits in degrees $\geq 0$
thereby proving the claim. This finishes the proof of (2).

\medskip\noindent
Assume $W = Y$ but make no assumptions on the algebraic space $Y$.
Since we have (2), we see from
Simplicial Spaces, Lemma \ref{spaces-simplicial-lemma-fppf-neg-ext-zero-hom}
that the functor $F$ given by $F(V) = \Hom_{\mathcal{O}_{X_V}}(K_V, E_V)$
is a sheaf\footnote{To check the sheaf property
for a covering $\{V_i \to V\}_{i \in I}$ first consider the
{\v C}ech fppf hypercovering $a : V_\bullet \to V$ with
$V_n = \coprod_{i_0 \ldots i_n} V_{i_0} \times_V \ldots \times_V V_{i_n}$
and then set $U_\bullet = V_\bullet \times_{a, V} X_V$. Then
$U_\bullet \to X_V$ is an fppf hypercovering to which we may
apply Simplicial Spaces, Lemma
\ref{spaces-simplicial-lemma-fppf-neg-ext-zero-hom}.}
on $(\textit{Spaces}/Y)_{fppf}$. Thus to prove that $F$
is an algebraic space and that $F \to Y$ is affine and of
finite presentation, we may work \'etale locally on $Y$; see
Bootstrap, Lemma \ref{bootstrap-lemma-locally-algebraic-space-finite-type}
and
Morphisms of Spaces, Lemmas \ref{spaces-morphisms-lemma-affine-local} and
\ref{spaces-morphisms-lemma-finite-presentation-local}. We conclude
that it suffices to prove $F$ is an affine algebraic space of
finite presentation over $Y$ when $Y$ is an affine scheme. In this
case we go back to our pseudo-coherent complex $L \in D(\mathcal{O}_Y)$.
Since $H^i(L) = 0$ for $i > 0$, we can represent $L$ by a complex
of the form
$$
\ldots \to \mathcal{O}_Y^{\oplus m_1} \to \mathcal{O}_Y^{\oplus m_0}
\to 0 \to \ldots
$$
with the last term in degree $0$, see More on Algebra, Lemma
\ref{more-algebra-lemma-pseudo-coherent}. Combining the two displayed formulas
earlier in the proof we find that
$$
F(V) =
\Ker(
\Hom_V(\mathcal{O}_V^{\oplus m_0}, \mathcal{O}_V)
\to 
\Hom_V(\mathcal{O}_V^{\oplus m_1}, \mathcal{O}_V)
)
$$
In other words, there is a fibre product diagram
$$
\xymatrix{
F \ar[d] \ar[r] & Y \ar[d]^0 \\
\mathbf{A}_Y^{m_0} \ar[r] & \mathbf{A}_Y^{m_1}
}
$$
which proves what we want.
\end{proof}

\end{document}
